\documentclass[11pt]{article}
\usepackage[utf8]{inputenc}
\usepackage[margin=2.5cm]{geometry}
\usepackage{amsmath}
\usepackage{amssymb}
\usepackage{booktabs}
\usepackage{siunitx}
\usepackage{colortbl}
\usepackage{xcolor}

\title{Multipole-Expansion Fit Model for the In-Plane Magnetic Field
       of a Kuo Quadrupole Magnetic Tweezer}
\author{}
\date{}

\begin{document}
\maketitle
\vspace{-2em}

\section*{Context}
This document specifies the forward model used to fit the magnetic
field in the workspace plane ($z = 0$) of the Kuo Quadrupole
magnetic tweezer (4 in-plane poles, fillet radius
\SI{45}{\micro\metre}, pole-tip distance
$R_{\text{pole}} = \SI{500}{\micro\metre}$). The model replaces the
classical point-charge ansatz by a component-wise real polynomial
expansion, which is the natural representation of the projection of
a 3D harmonic field onto the workspace plane.

% =====================================================================
\section{Forward model}
\label{sec:forward}

For each Cartesian field component
$\alpha \in \{x, y, z\}$, the model expresses
$B_\alpha$ at any planar point $(x, y)$ as a real polynomial in
the normalized coordinates $\tilde{x} \equiv x / r_0$ and
$\tilde{y} \equiv y / r_0$:
%
\begin{equation}
\label{eq:model}
\boxed{\;
B_\alpha(x, y) \;=\;
\sum_{\substack{i,\, j \,\ge\, 0 \\ i + j \,\le\, M}}
   c^{(\alpha)}_{ij}
   \left(\frac{x}{r_0}\right)^{i}
   \left(\frac{y}{r_0}\right)^{j}
\;,\qquad
\alpha \in \{x, y, z\}
\;}
\end{equation}
%
The three components $B_x$, $B_y$, $B_z$ are fit \emph{independently}
(no coupling is imposed across components).
The symbols are summarized in Table~\ref{tab:symbols}.

\begin{table}[htbp]
\centering
\caption{Symbol definitions for the forward model in
         Eq.~\eqref{eq:model}.}
\label{tab:symbols}
\begin{tabular}{c l l}
\toprule
Symbol & Meaning & Value / range \\
\midrule
$\alpha$           & Cartesian field component index & $\{x, y, z\}$ \\
$(x, y)$           & Field-evaluation point in workspace plane &
                     $z = 0$ slice \\
$r_0$              & Normalization length (conditioning) &
                     $\SI{500}{\micro\metre}$ ($= R_{\text{pole}}$) \\
$M$                & Polynomial truncation total degree &
                     $1, 2, \dots$ (e.g.\ 4, 6, 8, 12) \\
$(i, j)$           & Multi-index of basis monomial & $i + j \le M$ \\
$c^{(\alpha)}_{ij}$ & Real expansion coefficient (free parameter)
                     & determined by least squares \\
\bottomrule
\end{tabular}
\end{table}

The number of free coefficients per component is
%
\begin{equation}
\label{eq:dof}
K(M) \;=\; \binom{M + 2}{2} \;=\; \frac{(M+1)(M+2)}{2},
\end{equation}
%
so the total number of free parameters in the joint
$(B_x, B_y, B_z)$ model is $3\,K(M)$.

% =====================================================================
\section{Coefficient interpretation}
\label{sec:interpretation}

The polynomial basis
$\{\tilde{x}^{i} \tilde{y}^{j} : i + j \le M\}$
is a complete basis for analytic functions on a disk and contains
the conventional in-plane multipole expansion as a subset.
A monomial of total degree $n = i + j$ corresponds to multipole
moments of order $\le n$. Specifically, the change of basis to
$r^{n} \cos(m\phi)$ and $r^{n} \sin(m\phi)$
($r^{2} = \tilde{x}^{2} + \tilde{y}^{2}$,
$\phi = \arctan(y/x)$, $0 \le m \le n$, with $n - m$ even)
gives the standard multipole interpretation listed in
Table~\ref{tab:interp}.

\begin{table}[htbp]
\centering
\caption{Physical meaning of the coefficients
         $c^{(\alpha)}_{ij}$ grouped by total degree
         $n = i + j$, expressed for each component
         $\alpha \in \{x, y, z\}$. ``Multipole order'' refers to the
         angular index $m$ in the $\cos(m\phi)$ / $\sin(m\phi)$
         basis. Each row contributes $(n + 1)$ monomials per component.}
\label{tab:interp}
\begin{tabular}{c c c l}
\toprule
$n = i+j$ & Coefficients & Multipole order(s) $m$ & Physical meaning \\
\midrule
0 & $c^{(\alpha)}_{00}$
  & $0$
  & Uniform background (offset, measurement bias) \\
1 & $c^{(\alpha)}_{10},\; c^{(\alpha)}_{01}$
  & $1$
  & Linear field gradient (dipole-equivalent) \\
2 & $c^{(\alpha)}_{20},\; c^{(\alpha)}_{11},\; c^{(\alpha)}_{02}$
  & $0,\; 2$
  & \cellcolor{yellow!25}\textbf{Quadrupole gradient (dominant mode for 4 in-plane poles)} \\
3 & $c^{(\alpha)}_{30},\, c^{(\alpha)}_{21},\, c^{(\alpha)}_{12},\, c^{(\alpha)}_{03}$
  & $1,\; 3$
  & Sextupole correction \\
4 & $c^{(\alpha)}_{40} \dots c^{(\alpha)}_{04}$
  & $0,\; 2,\; 4$
  & Octupole correction \\
$\vdots$ & $\vdots$ & $\vdots$ & higher-order corrections \\
$M$ & $c^{(\alpha)}_{M,0} \dots c^{(\alpha)}_{0,M}$
    & up to $m = M$
    & truncation level chosen by accuracy requirement \\
\bottomrule
\end{tabular}
\end{table}

\paragraph{Relation to analytic 2D Taylor expansion.}
A pure 2D analytic Taylor expansion of the complex field would write
$B_x - i B_y = \sum_{m=0}^{M} a_m\, z^{m}$
with $z = x + i y$, contributing only $2(M+1)$ real DOF per component
pair. The polynomial basis in Eq.~\eqref{eq:model} is a
strict superset, containing additional non-analytic monomials such
as $r^{2}\, x^{i} y^{j}$ that the analytic ansatz cannot represent.
Section~\ref{sec:why} explains why these extra terms are required.

% =====================================================================
\section{Why a component-wise polynomial basis}
\label{sec:why}

Inside the air region surrounding the workspace, the magnetic field
$\mathbf{B}$ satisfies the source-free 3D equations
$\nabla \cdot \mathbf{B} = 0$ and
$\nabla \times \mathbf{B} = \mathbf{0}$,
so each Cartesian component is a 3D harmonic function.
Restricted to the workspace plane $z = 0$, however, the 2D divergence
of the in-plane components is in general \emph{not} zero:
%
\begin{equation}
\label{eq:divxy}
\left.\frac{\partial B_x}{\partial x}
   + \frac{\partial B_y}{\partial y}\right|_{z=0}
\;=\; -\,\left.\frac{\partial B_z}{\partial z}\right|_{z=0}
\;\neq\; 0.
\end{equation}
%
For a $z$-symmetric pole geometry the right-hand side is even in $z$
and finite at $z = 0$, even though $B_z(z = 0) = 0$ by symmetry.
Consequently, the in-plane field components do not satisfy the 2D
Cauchy--Riemann conditions that an analytic Taylor expansion in
$z = x + i y$ would require. The component-wise real polynomial
basis in Eq.~\eqref{eq:model} does \emph{not} impose such a
constraint and therefore captures the full projection of the 3D
harmonic field onto the workspace plane, including the non-analytic
$r^{2k}$ corrections.

% =====================================================================
\section{Fit procedure}
\label{sec:fit}

Given $N$ FEM samples
$\{(x_k, y_k,\, B_x^{(k)},\, B_y^{(k)},\, B_z^{(k)})\}_{k=1}^{N}$
(typically restricted to $z = 0 \pm z_{\text{tol}}$ and to a sampling
disk $x_k^{2} + y_k^{2} \le R^{2}$, with iron-pole wedges removed),
construct the design matrix
%
\begin{equation}
A_{k,\,p(i,j)}
\;=\; \left(\frac{x_k}{r_0}\right)^{i}
      \left(\frac{y_k}{r_0}\right)^{j},
\qquad i + j \le M,
\end{equation}
%
where $p(i,j)$ enumerates the $K(M)$ basis monomials. The $A$ matrix
has shape $N \times K$ and is shared across components. Each
component is then fit by ordinary linear least squares,
%
\begin{equation}
\label{eq:lstsq}
\hat{\mathbf{c}}^{(\alpha)}
\;=\; \arg\min_{\mathbf{c}^{(\alpha)} \in \mathbb{R}^{K}}
      \big\|\, A\, \mathbf{c}^{(\alpha)} - \mathbf{B}^{(\alpha)} \,\big\|_{2}^{\,2},
\qquad \alpha \in \{x, y, z\},
\end{equation}
%
which has a closed-form normal-equation solution and contains no
nonlinear parameters and no iterative optimization.

% =====================================================================
\section{Degrees of freedom}
\label{sec:dof}

\begin{table}[htbp]
\centering
\caption{Free parameters per component, $K(M) = (M+1)(M+2)/2$, and
         total parameters $3K(M)$ for the joint
         $(B_x, B_y, B_z)$ model. The rightmost column lists the
         highest multipole order $m_{\max} = M$ that the
         truncation can represent.}
\label{tab:dof}
\begin{tabular}{c c c c}
\toprule
$M$ & $K(M)$ per component & $3K(M)$ total &
      Highest multipole $m_{\max}$ \\
\midrule
1  &  3  &   9 & dipole \\
2  &  6  &  18 & quadrupole \\
3  & 10  &  30 & sextupole \\
4  & 15  &  45 & octupole \\
6  & 28  &  84 & dodecapole \\
8  & 45  & 135 & 16-pole \\
10 & 66  & 198 & 20-pole \\
12 & 91  & 273 & 24-pole \\
\bottomrule
\end{tabular}
\end{table}

% =====================================================================
\appendix
\section{Explicit basis listing for $M = 4$}
\label{app:m4}

For concreteness, the $K(4) = 15$ basis monomials of total degree
$\le 4$ used in Eq.~\eqref{eq:model} are
%
\begin{align*}
n = 0:\quad & 1 \\
n = 1:\quad & \tilde{x},\; \tilde{y} \\
n = 2:\quad & \tilde{x}^{2},\; \tilde{x}\tilde{y},\; \tilde{y}^{2} \\
n = 3:\quad & \tilde{x}^{3},\; \tilde{x}^{2}\tilde{y},\;
              \tilde{x}\tilde{y}^{2},\; \tilde{y}^{3} \\
n = 4:\quad & \tilde{x}^{4},\; \tilde{x}^{3}\tilde{y},\;
              \tilde{x}^{2}\tilde{y}^{2},\;
              \tilde{x}\tilde{y}^{3},\; \tilde{y}^{4}
\end{align*}
%
with $\tilde{x} = x / r_0$, $\tilde{y} = y / r_0$.
Each of $B_x, B_y, B_z$ uses these same 15 monomials, fit
independently, for a total of $3 \times 15 = 45$ free parameters
at $M = 4$.

\end{document}
